% Pages/introducao.tex

\lettrine[lines=3, lhang=0.1, findent=2pt]{\ringmfamily O}{ Sistema Overgrown} é uma proposta original concebida para proporcionar campanhas medievais divertidas, priorizando a narrativa e a experiência de jogo e removendo complexidade técnica desnecessária de outros sistemas. Aqui, a burocracia de outros sistemas cede lugar à fluidez: o personagem é aquilo que deseja ser. Suas habilidades e talentos não são limitados por classes rígidas, mas moldados por sua ligação direta com entidades superiores e pela lógica intrínseca deste mundo.

\section*{O Mundo Fragmentado}
O universo de Overgrown é vasto e fragmentado, composto por continentes independentes com culturas e políticas próprias. No coração desta realidade reside uma ilha singular: \textbf{A Guilda}. Desde o primeiro grande Cataclismo, este local tornou-se o refúgio dos deuses, servindo como o eixo de equilíbrio, de controle, e manutenção de toda a existência.

\section*{O Primeiro e Segundo Cataclismo}
O ano de 2020 marcou o início do caos absoluto. Diante da iminente ruína de sua criação por guerras e caos absoluto, os Celestiais desceram ao plano terrestre para salvar a Terra do Primeiro Cataclismo. Contudo, a intervenção divina trouxe consequências imprevistas. 

Onde um ser superior tocava o solo, abria-se uma fenda de energia primordial. Essa energia era capaz de alterar a matéria e todas as leis físicas e naturais. Esta energia espalhou-se como uma praga, transformando tudo o que tocava e permitindo que seres celestiais se manifestassem fisicamente. 

Este evento foi batizado de \textbf{Overgrown}, dando inicio ao Segundo Cataclismo. O mundo rompeu-se em terremotos e maremotos, devastando nações inteiras. Alguns deuses se espalharam pelo mundo para ajudar na contenção da destruição causada, enquanto outros deuses recuaram para o centro do mundo, levando com eles grandes quantidades de pessoas, animais e criaturas, onde foi criada a Guilda, como medida segurança como o \textbf{Ultimo Bastião da Civilização} em caso de um apocalipse total.


\section*{A Era da Expansão}
Por 2.000 anos, a Guilda permaneceu isolada, acreditando realmente ser o \textbf{Ultimo Bastião da Civilização}. No entanto, em 4020, as fronteiras foram abertas, dando início à \textbf{Grande Expansão}. Pelos 500 anos seguintes, alianças foram formadas, bases foram estabelecidas e aventureiros passaram a ser treinados dentro e fora da Guilda. Foi criado um sistema de missões da Guilda, para auxiliar e desenvolver todos os países.

\begin{rpgboxtips}{Missões da Guilda}
A Guilda é uma organização de postura neutra. Ela não impõe julgamentos morais sobre o que é certo ou errado nas missões realizadas fora de seus territórios, reconhecendo que cada país e cada continente possuem sua própria cultura, leis e valores. Em geral, missões iniciadas dentro dos domínios da Guilda tendem a ser mais controladas e pacíficas, enquanto aquelas conduzidas em territórios estrangeiros refletem as tensões, conflitos e instabilidades próprias de cada região.
\end{rpgboxtips}

\originlink{res:null}
\section*{O Surgimento do Null}
A paz longa e duradora foi estilhaçada em 4521. A cidade capital dos anões, Yroln, desapareceu subitamente, como se jamais tivesse existido. Coincidindo com esse evento, o surgimento do \textbf{Null} — a Energia Impossível, ou apenas \textbf{Energia}. Esta energia se espalhou como vírus desconhecido e começou a consumir a própria mana, espalhando-se com uma instabilidade superior a qualquer poder conhecido. Deuses que não se adaptaram pereceram em liberações violentas de energia, enquanto os sobreviventes tornaram-se ainda mais poderosos.

\section*{O Selo de Divindade}
Em 4522, os antigos \textit{Oculys} foram abolidos. No seu lugar, instituiu-se o \textbf{Sistema de Selos de Divindade}. Estas marcas mágicas, gravadas nas costas dos aventureiros, servem como um elo direto com o \textbf{Pináculo} — uma torre inacessível que abriga o conhecimento acumulado dos celestiais. Todo selo existente é igual e representa o vínculo individual com o Pináculo e seu aventureiro.

\begin{rpgboxtips}{O Pináculo}
    É possível se conectar ao Pináculo sem fazer parte da Guilda, por meio de magia, no entanto, o custo para criar um selo é muito alto para pessoas comuns e conectar alguém sem permissão da Guilda é considerado uma prática proibida e criminosa. 
\end{rpgboxtips}

